\section{The geometry is derived, not assumed}\label{sec:geometry}

Let $\PP(x)$ be a $\C^1$ family of orthogonal projectors of constant rank $k$, parametrised by
$x=(x^1,\dots,x^m)$ in an open set. Write $\partial_\mu=\partial/\partial x^\mu$.

\begin{definition}[Quantum geometric tensor of a projector family]\label{def:qgt}
\begin{equation}\label{eq:qgt}
   Q_{\mu\nu}(x):=\Tr\bigl(\PP\,\partial_\mu\PP\,\partial_\nu\PP\bigr),\qquad
   g_{\mu\nu}:=\operatorname{Re}Q_{\mu\nu},\qquad
   \omega_{\mu\nu}:=-2\operatorname{Im}Q_{\mu\nu}.
\end{equation}
\end{definition}

The placement of $\PP$ matters and is the point of the next remark.

\begin{theorem}[$Q$ is a Hermitian tensor and $g$ is a metric]\label{thm:qgt}
$Q$ is Hermitian, $g$ is symmetric and positive semidefinite, and for a single parameter with
variation $\dot\PP$,
\begin{equation}\label{eq:metric-norm}
   g(\dot\PP,\dot\PP)=\norm{\PP\dot\PP\QQ}_F^2=\tfrac12\norm{\dot\PP}_F^2=\tfrac12\Tr(\dot\PP^2).
\end{equation}
The identity $g_\PP(V,V)=\tfrac12\Tr(V^2)$ is the June paper's \Cc{} \cite{corpus}, where
the metric is defined that way; what is added here is the derivation, which makes it the
special case of the redistribution theorem: $\dot\PP=\PP\dot\PP\QQ+\QQ\dot\PP\PP$
(differentiate $\PP^2=\PP$), so Theorem~\ref{thm:norm} applied to $\Omega=\dot\PP$ gives
$\norm{\PP\dot\PP\QQ}_F^2=\tfrac12\norm{\dot\PP}_F^2$. \T\ for the derivation, \Cc\ for the
identity; the kernel checks both halves (\lean{redistribution\_norm},
\lean{half\_commutator\_norm}).
\end{theorem}

\begin{proof}
Hermiticity: $\overline{Q_{\nu\mu}}=\overline{\Tr(\PP\partial_\nu\PP\partial_\mu\PP)}=\Tr\bigl((\PP\partial_\nu\PP\partial_\mu\PP)^*\bigr)=\Tr(\PP\partial_\mu\PP\partial_\nu\PP)=Q_{\mu\nu}$.
Positivity: put $Z_\mu:=\QQ\partial_\mu\PP\,\PP$. Since $\partial_\mu\PP=\PP\partial_\mu\PP\QQ+\QQ\partial_\mu\PP\PP$,
$\Tr(\PP\partial_\mu\PP\partial_\nu\PP)=\Tr(\PP\partial_\mu\PP\,\QQ\,\partial_\nu\PP)=\Tr(Z_\mu^*Z_\nu)$,
so $Q_{\mu\nu}=\ip{Z_\mu}{Z_\nu}_F$ and $g_{\mu\mu}=\norm{Z_\mu}_F^2\ge0$. Equation
\eqref{eq:metric-norm} is the special case $\mu=\nu$ together with Theorem~\ref{thm:norm}.
\end{proof}

\begin{remark}[A second correction: the tensor must carry the projector]\label{rem:qgt-correction}
An earlier formulation used $Q^\circ_{\mu\nu}=\Tr(\partial_\mu\PP\,\partial_\nu\PP)$ as the
quantum geometric tensor. That object is real: $\partial_\mu\PP$ is self-adjoint, so
$(\partial_\mu\PP\partial_\nu\PP)^*=\partial_\nu\PP\partial_\mu\PP$ and $\Tr$ of the product is
real, hence $\operatorname{Im}Q^\circ=0$ identically and $Q^\circ$ cannot carry curvature. What
it is, by Theorem~\ref{thm:qgt}, is twice the metric: $Q^\circ=2g$. The corrected tensor $Q$ of
\eqref{eq:qgt} has a genuine imaginary part, and it is the corpus's own form: its later paper
\Cc{} \cite{iqgt} writes the intrinsic tensor as $Q_{\PP}(V,W)=\Tr(\PP VW)$, which is
\eqref{eq:qgt} for $V=\partial_\mu\PP$, $W=\partial_\nu\PP$; the object $Q^\circ$ was not in fact
a form the corpus ever used --- the June paper states $Q_\PP(V,W)=\Tr(\PP VW)$ already --- so
this remark corrects an earlier draft of \emph{this} paper, not the corpus. \V\ Both facts are measured in
\gate{results/swarm/check\_D2\_m7.py} (literal imaginary part exactly zero; nonzero curvature
$0.358678$): two independent reports found the same error, which is why it is recorded here
rather than quietly repaired.
\end{remark}

\begin{theorem}[Curvature and its periods]\label{thm:curv}
For a $\C^2$ family, $\omega=\tfrac12\omega_{\mu\nu}\df x^\mu\wedge\df x^\nu$ is closed,
$\df\omega=0$. For a closed oriented surface $\shape$ and rank $k=1$,
$\frac1{2\pi}\int_\shape\omega\in\mathbb{Z}$; for higher rank the trace of the curvature is
likewise integral (Chern--Weil). \T\ closedness, from $\PP^2=\PP$ and the symmetry of mixed
partials; the integrality is classical \Cc{} (Chern--Weil theory; for the rank-one case the
Berry phase about a cone is $\pi$ per unit winding, \S\ref{sec:hierarchy}).
\end{theorem}

\begin{theorem}[Metric--curvature bound]\label{thm:bound}
For every rank $k$ and every pair of tangent directions $V,W$,
\begin{equation}\label{eq:bound}
   \abs{\omega(V,W)}\;\le\;2\sqrt{g(V,V)g(W,W)-g(V,W)^{2}}\;\le\;2\sqrt{g(V,V)g(W,W)} ,
   \qquad\text{and}\qquad \abs{\operatorname{Im}Q(V,W)}\le\tfrac12\abs{\omega(V,W)} .
\end{equation}
All three are Gram-determinant inequalities for the pair $Z_V,Z_W$ of
Theorem~\ref{thm:qgt}: $\abs{Q(V,W)}^2=g(V,W)^2+\tfrac14\omega(V,W)^2$ and
$g(V,V)g(W,W)\ge\abs{Q(V,W)}^2$. The first is sharp, with equality exactly when $V,W$ are
linearly dependent over $\C$; the second drops the Gram term and is sharp only when $g(V,W)=0$.
\T\ (Cauchy--Schwarz and the Gram determinant). Sharpness: \V\
\gate{results/swarm/check\_D3\_m1.py}.
\end{theorem}

\begin{remark}[The sharp form is the corpus's; an earlier draft of this paper was weaker]
\label{rem:factor2}
Theorem 12 of the June paper \Cc{} \cite{corpus} states
$\abs{\Omega_\PP(V,W)}\le2\sqrt{g_\PP(V,V)g_\PP(W,W)-g_\PP(V,W)^2}$, which with
$\Omega_\PP=2\operatorname{Im}Q=\omega$ is the first inequality of \eqref{eq:bound} --- the sharp
Gram-determinant form, valid for every rank because a Gram determinant does not see the rank.
An earlier draft of this paper offered $\abs{\omega}\le2\sqrt{g(V,V)g(W,W)}$ as a correction of
that bound. It is not one: the second inequality of \eqref{eq:bound} is the first with the Gram
term discarded, so the corpus's statement is the stronger of the two and this paper's was the
weaker. What survives of the earlier point is only that the constant $2$ belongs to the
\emph{curvature} and the constant $1$ to $\operatorname{Im}Q$ in the normalisation of
\eqref{eq:qgt} where $\omega=-2\operatorname{Im}Q$; in the corpus's normalisation
$Q=g+\tfrac{i}{2}\Omega$ the same statement reads $\abs{\Omega}\le2\sqrt{\cdot}$ and
$\abs{\operatorname{Im}Q}\le\sqrt{\cdot}$, which is the same pair of constants.
\end{remark}


\begin{remark}[What the section says, in one sentence]
The metric and the curvature are the symmetric and antisymmetric parts of one derived object,
the tensor $Q$, and the metric is not an input but a squared norm of a block of $\dot\PP$ ---
which is why \S\ref{sec:redistribution} had to come first.
\end{remark}
